\maketitle

\begin{tcolorbox}[title={Main endpoint theorem (Fefferman (B), periodic case)}, colback=blue!2!white, colframe=blue!55!black]
\textbf{Classical endpoint.}
Let $\nu>0$ and let $u_0\in C^\infty(\T^3;\R^3)$ satisfy $\nabla\cdot u_0=0$.
Then there exists a unique smooth periodic Navier--Stokes solution pair $(u,p)$ on
$\T^3\times[0,\infty)$ satisfying
\[
\partial_tu+(u\cdot\nabla)u+\nabla p=\nu\Delta u,
\qquad \nabla\cdot u=0,
\qquad u|_{t=0}=u_0,
\]
with the fixed mean-zero pressure gauge. This statement is obtained first as
Corollary~\ref{cor:official-periodic-structural-corollary} and then restated in the official classical form as
Theorem~\ref{thm:official-periodic-pde-theorem}.
\end{tcolorbox}

\begin{tcolorbox}[breakable, title={Proof-class, claim-class, and source-use locks}, colback=yellow!3!white, colframe=yellow!45!black]
\textbf{Proof class.} This manuscript is a mathematical non-PDE-native constraint-formalism impossibility proof of Fefferman's periodic Navier--Stokes statement~(B). It is not a continuation-estimate proof, not a constructive solution-formula proof, and not a claim that admissibility supplies a missing vortex-stretching bound.

\medskip
\textbf{Claim class.} The proof fixes the official periodic problem as a non-degenerate same-scope proof/refutation regime. It then shows that a hypothetical finite-horizon breakdown induces a same-scope endpoint condition, that such a condition has only the endpoint statuses classified below, and that no licit same-scope finite-horizon counterexample completes. The conclusion is obtained by impossibility of counterexample completion, not by an analytic continuation template.

\medskip
\textbf{Clay/Fefferman method-neutrality.} The official Clay rules make eligibility depend on a complete mathematical solution that satisfactorily answers the questions in the official problem description after publication, scrutiny, and general acceptance; they do not prescribe a PDE-native continuation-estimate template, a vortex-stretching bound, a weak-to-strong regularity route, or a constructive solution formula as the required method class~\cite{cmi-millennium-rules}. Fefferman's official Navier--Stokes description likewise states that solvers are given reasonable leeway while retaining the heart of the problem, and it asks for proof of one of four official alternatives, including periodic existence and smoothness statement~(B)~\cite{fefferman-navierstokes-problem}. Thus a PDE-native method is not a Clay-rule prerequisite. What is required here is still the full mathematical endpoint: a proof of the official periodic statement itself.

\medskip
\textbf{Source-use lock.} The live proof chain is restricted to theorem-local material in the manuscript, Appendix~A, Appendix~B, Appendix~E for the analytic atlas used by Part~I, and public official-problem/source-alignment material where cited. Private notes, companion slogans, validation packets, and non-load-bearing audit files do not carry the endpoint. Appendix~\ref{sec:internal-ledger} records the source-governance and Lean-support boundary.
\end{tcolorbox}

\begin{tcolorbox}[breakable, title={Clay rules do not require a PDE-native proof method}, colback=yellow!2!white, colframe=orange!65!black]
The phrase ``Navier--Stokes problem'' names the official PDE endpoint, but the Clay rules do not require a proof to be PDE-native in method. They require a complete mathematical solution to the Problem as defined in its official description, publication in a qualifying outlet, at least two years of scrutiny, general acceptance in the global mathematics community, and CMI examination~\cite{cmi-millennium-rules}. Fefferman's official description is also method-neutral in this sense: it grants reasonable leeway to solvers while retaining the heart of the problem and asks for proof of one of the four stated alternatives~\cite{fefferman-navierstokes-problem}.

\medskip
Accordingly, a referee objection of the form ``this is not a PDE continuation estimate'' identifies the proof class but does not by itself refute the manuscript. A same-mode objection must show that the non-PDE-native constraint-formalism route fails to prove the official periodic statement (B), fails to preserve Fefferman's data/equation/gauge/continuation target, or fails at one of the theorem-local links in the declared proof spine. This paragraph does not relax the endpoint: the theorem still proves the classical periodic PDE statement stated above.
\end{tcolorbox}

\begin{tcolorbox}[breakable, title={Counterexamplehood lock}, colback=cyan!2!white, colframe=cyan!55!black]
A finite-time expression such as $T_*(u_0)<\infty$ is not by itself a theorem-bearing refutation of Fefferman~(B). To refute the official universal statement it must be admitted in the ambient proof/refutation regime as an official datum, an official realized maximal classical lineage, a finite-horizon condition, and a same-scope non-continuation claim. Ordinary ZFC or classical-analysis counterexamplehood therefore already requires admissible formation, proof-standing, fixed reference, same-target fidelity, and final theorem consequence. Those are precisely the kernel roles at theorem-bearing use level. If an alleged object lacks those roles, it is not a counterexample in the ambient proof system either.

\medskip
The relevant theorem-local nodes are Theorem~\ref{thm:ambient-proof-system-subsumption}, Lemma~\ref{lem:counterexamplehood-theorem-bearing-construction}, Theorem~\ref{thm:fefferman-counterexample-subsumption}, and Corollary~\ref{cor:no-ambient-proof-system-escape}.
\end{tcolorbox}

\begin{tcolorbox}[breakable, title={Read this first: fixed-scope horizon proof spine}, colback=blue!2!white, colframe=blue!60!black]
The live proof route is the following single spine:
\[
\begin{aligned}
&\text{official periodic proof/refutation regime}
\Rightarrow K^{NS}_{\mathrm{per}}
\Rightarrow \mathcal I_{\mathrm{car}}(\T^3,\nu)\text{ unique} \\
&\Rightarrow \mathcal R_{NS}(x)\text{ is one same-scope realized lineage}
\Rightarrow [T_*(x)<\infty]_i
\Rightarrow B^{NS}_{\mathrm{hor}}(x,T_*(x)) \\
&\Rightarrow \text{endpoint-status trichotomy}
\Rightarrow \text{non-witness branches collapse} \\
&\Rightarrow \text{standing-positive witness branch is exhausted} \\
&\Rightarrow T_*(x)=\infty
\Rightarrow \text{Fefferman (B), periodic case.}
\end{aligned}
\]
The notation $[T_*(x)<\infty]_i$ denotes the local exact-complement countercase opened in the endpoint proof. It is not an additional object-level premise. The horizon bridge object $B^{NS}_{\mathrm{hor}}(x,T_*(x))$ packages only the finite-horizon countercase data: fixed carrier, realized-lineage identity, pre-horizon witness material, horizon discriminator, metric-bookkeeping discipline, and branch extraction. It does not contain the conclusion $T_*(x)=\infty$.

\medskip
The detailed theorem ladder is recorded in Appendix~\ref{sec:ns-theorem-ladder}. The route-locus audit for analytic and endpoint alternatives is recorded in Appendix~\ref{sec:ns-route-locus-audit}.
\end{tcolorbox}

\begin{tcolorbox}[breakable, title={No one-step counterexample promotion}, colback=cyan!1!white, colframe=cyan!45!black]
The proof does not use the inference
\[
T_*(x)<\infty\quad\Longrightarrow\quad\text{completed counterexample to Fefferman (B)}.
\]
The route is instead
\[
[T_*(x)<\infty]_i
\Rightarrow \text{finite-horizon endpoint condition}
\Rightarrow B^{NS}_{\mathrm{hor}}(x,T_*)
\Rightarrow \text{trichotomy}
\Rightarrow \bot.
\]
The contradiction discharges the local finite-horizon assumption itself. It does not discharge a strengthened premise such as ``finite-time breakdown plus official endpoint-governed negative status.''
\end{tcolorbox}

\begin{tcolorbox}[breakable, title={Vortex-stretching seam discipline}, colback=green!2!white, colframe=green!45!black]
The vortex-stretching channel, BKM-type quantities, critical norms, metric divergences, and profile/concentration diagnostics remain valid analytic diagnostics and obstruction witnesses. The manuscript does not bound them by a new estimate. Their endpoint role is narrower: they do not become independent same-scope endpoint authority. On the live route they are licensed metric bookkeeping or quotient-visible witness bookkeeping subordinate to the finite-horizon endpoint condition and the carrier-side horizon discriminator. If promoted as independent continuation authority, they become an illicit same-scope metric gate or a scope move.

\medskip
This seam is localized in Proposition~\ref{prop:metric-divergence-skin-non-authority}, Proposition~\ref{prop:quotient-visible-witness-vs-illicit-metric-gate}, Theorem~\ref{thm:metric-non-authority-periodic-carrier}, Theorem~\ref{thm:dynamic-stretching-compatibility}, and Appendix~\ref{sec:ns-route-locus-audit}.
\end{tcolorbox}

\begin{tcolorbox}[breakable, title={Fast denial targets}, colback=gray!3!white, colframe=gray!65!black]
\small
\renewcommand{\arraystretch}{1.12}
\begin{tabularx}{0.98\textwidth}{@{}L{0.42\textwidth}Y@{}}
\toprule
\textbf{Objection} & \textbf{Theorem-local denial target}\\
\midrule
No PDE-native estimate is supplied. & Clay/Fefferman method-neutrality box; proof-class lock; vortex-stretching seam discipline; Theorem~\ref{thm:counterexample-soundness-fixed-same-scope}.\\
A classical counterexample escapes the kernel. & Theorem~\ref{thm:ambient-proof-system-subsumption}; Theorem~\ref{thm:fefferman-counterexample-subsumption}.\\
Finite $T_*$ automatically refutes Fefferman~(B). & Theorem~\ref{thm:finite-horizon-failure-induces-endpoint-condition}; Theorem~\ref{thm:counterexample-soundness-fixed-same-scope}; horizon reductio in Section~\ref{subsec:same-scope-continuation-failure-dichotomy}.\\
The realized lineage can be admissible before $T_*$ and non-admissible at $T_*$. & Theorem~\ref{thm:carrier-to-lineage-binding}; Theorem~\ref{thm:no-admissibility-dynamics-realized-lineages}; Theorem~\ref{thm:st4-standing-conservation-realized-lineages}.\\
Metric divergence carries endpoint authority. & Proposition~\ref{prop:metric-divergence-skin-non-authority}; Proposition~\ref{prop:quotient-visible-witness-vs-illicit-metric-gate}; Theorem~\ref{thm:metric-non-authority-periodic-carrier}.\\
A weak/profile/alternative formulation is a counterexample. & Fixed-domain comparison class, Definition~\ref{def:fixed-domain-comparison-class}; route-locus L8 and L14 in Appendix~\ref{sec:ns-route-locus-audit}.\\
Part III reopens primitive admissibility. & Corollary~\ref{cor:no-circular-burden-after-carrier-packet}; Proposition~\ref{prop:partIII-transport-only}; Proposition~\ref{prop:partIII-same-continuation-notion}.\\
The endpoint does not restate Fefferman~(B). & Theorem~\ref{thm:primitive-to-classical-transfer-cascade}; Corollary~\ref{cor:official-periodic-structural-corollary}; Theorem~\ref{thm:official-periodic-pde-theorem}.\\
\bottomrule
\end{tabularx}
\end{tcolorbox}

\begin{abstract}
This manuscript proves the periodic three-dimensional Navier--Stokes endpoint in Fefferman's statement~(B) by a fixed-scope constraint-formalism impossibility route. The argument does not add a vortex-stretching estimate or a new continuation criterion. It proves that the official periodic proof/refutation regime instantiates the fixed-scope kernel, identifies the standard periodic carrier with the unique admissible interior, realizes every official smooth divergence-free datum as one same-scope maximal classical lineage, and then treats a hypothetical finite horizon as a local exact-complement countercase. The horizon bridge object associated to that countercase is exhausted by a three-way endpoint-status classification: standing-positive witness, terminal standing failure, or illicit same-scope gate/scope move. The non-witness branches collapse locally, and the remaining witness branch is exhausted and eliminated by the branch-forcing packet. Hence no finite-horizon same-scope counterexample completes, so every official periodic smooth datum has a unique global smooth periodic solution with the fixed mean-zero pressure gauge.
\end{abstract}

\tableofcontents
\clearpage
